% =====================================================================
%  How the second slit of towards_double_slit switches on
%  Build:  pdflatex WaveVisualizationModifier.tex   (run twice for the ToC)
% =====================================================================
\documentclass[11pt,a4paper]{article}

\usepackage[T1]{fontenc}
\usepackage[utf8]{inputenc}
\usepackage[margin=2.2cm]{geometry}
\usepackage{lmodern}
\usepackage{microtype}
\usepackage{amsmath,amssymb}
\usepackage{booktabs}
\usepackage{enumitem}
\usepackage{xcolor}
\usepackage{hyperref}

\hypersetup{
    colorlinks=true,
    linkcolor=blue!55!black,
    urlcolor=blue!55!black,
    citecolor=blue!55!black,
    pdftitle={WaveVisualizationModifier: a late source},
    pdfauthor={BlenderMathAnim}
}

\newcommand{\node}[1]{\texttt{#1}}
\newcommand{\clamp}{\operatorname{clamp}}

\title{A Slit That Opens Late\\
       \large How \node{WaveVisualizationModifier} lets an interference
       pattern grow instead of appearing at once}
\author{BlenderMathAnim --- \node{geometry\_nodes/modifier\_video\_interferences.py}}
\date{}

\begin{document}
\maketitle
\tableofcontents
\bigskip

% =====================================================================
\section{The question}
% =====================================================================

In \node{towards\_double\_slit} (\node{video\_interferences/scene\_interference.py})
a single slit has been radiating for the whole of the previous scene. At
$t_\text{open}=8\,$s the outer piece of the wall slides aside and a second
slit opens. What the surface behind the wall then shows is \emph{not} the
finished two-slit pattern. Instead a circular disturbance leaves the new
slit, and only the region it has already swept over carries the hyperbolic
nodal lines of the interference pattern; outside it the old single-slit
wave is untouched. After about seven seconds the disturbance has crossed the
whole surface and the stationary pattern is complete.

Every wave this modifier draws is a \emph{steady state}: the field of a
source that has been ringing for ever. Such a field is non-zero everywhere at
every instant, so a steady-state source that is simply switched on would make
the full pattern pop into existence in a single frame. This note explains the
few nodes that turn ``switched on'' into ``spreading out''.

% =====================================================================
\section{The steady state}
% =====================================================================

The surface is a grid in the $x$-$y$ plane whose vertices are lifted by
\begin{equation}
  u(\mathbf r,t)=\sum_j w_j(\mathbf r,t),\qquad
  w_j = A'\,\Big[J_0(k r_j)\cos\omega t+Y_0(k r_j)\sin\omega t\Big]
      = A'\,\mathrm{Re}\!\left[H^{(1)}_0(k r_j)\,e^{-i\omega t}\right],
  \label{eq:steady}
\end{equation}
with $r_j=|\mathbf r-\mathbf c_j|$ (held above \node{source\_radius} to stay
clear of the logarithmic pole of $Y_0$), $k=2\pi/\lambda$, $\omega=2\pi f$ and
$A'=A\pi/\sqrt\lambda$. Each $w_j$ is the exact outgoing solution of the 2+1
dimensional wave equation for a point source at $\mathbf c_j$, and because the
equation is linear their sum is a solution too: with two sources it is the
two-slit pattern.

The scene builds the surface behind the wall as
\begin{verbatim}
wave2 = WaveVisualizationModifier(name="DoubleSlitWave", size=8.0,
            sources=((-4, 0), (-4, -a)),
            source_impacts=(None, t_open + 0.5),
            wavelength=1.5, frequency=1.0, amplitude=0.31,
            material=get_texture("function", alpha_intensity=0.9),
            source_radius=2 / 10, time_shift=time_shift,
            intensity=True, intensity_gain=2.0)
\end{verbatim}
The grid has side 8 and is moved by $+4$ in $x$, so its left edge, local
$x=-4$, lies on the wall. Source~0 is the old slit (it slides from $y=0$ to
$y=+a$ before the second one opens); source~1 is the new slit at $y=-a$,
$a=1.25$. The one argument that matters here is
\node{source\_impacts}: \node{None} for the old slit, which rings for ever,
and the time $t_1=t_\text{open}+0.5=8.5\,$s for the new one --- the middle of
the one second during which the wall slides open.

% =====================================================================
\section{The causal envelope}
% =====================================================================

A source with an impact time $t_j$ is multiplied by an envelope that is
zero ahead of a circle expanding from $\mathbf c_j$ at the phase speed
$c=\lambda f$, and one behind it:
\begin{equation}
  \mathrm{env}_j(\mathbf r,t)=S\!\left(\frac{c\,(t-t_j)-r_j}{\Delta}\right),
  \qquad S(x)=\hat x^{2}(3-2\hat x),\quad \hat x=\clamp_{[0,1]}(x).
  \label{eq:env}
\end{equation}
\begin{itemize}[leftmargin=*]
  \item Before the impact $c(t-t_j)<0$, so $\hat x=0$ at every vertex and the
        source contributes nothing. The clamp is the whole reason the new slit
        is silent until it opens.
  \item After the impact the \emph{reach} $c(t-t_j)$ grows linearly. A vertex
        at distance $r_j$ starts to feel the source at $t=t_j+r_j/c$ and has
        the full steady-state contribution a time $\Delta/c$ later.
  \item $\Delta$ is the \node{FrontWidth} dial, one wavelength by default.
        The smoothstep $S$ has zero slope at both ends, so the edge of the
        disturbance does not read as a crease in the surface.
\end{itemize}

The surface therefore shows
\[
  u = w_0 + \mathrm{env}_1\,w_1 =
  \begin{cases}
    w_0 & r_1 > c(t-t_1) \qquad\text{(single slit, untouched)}\\
    w_0 + w_1 & r_1 < c(t-t_1)-\Delta \qquad\text{(full two-slit pattern)}
  \end{cases}
\]
with a band of width $\Delta$ in between. The disk inside which the pattern is
already stationary grows at speed $c$ around the new slit. That is the
``slow growth'' seen in the shot: the interference is not faded in, it is
\emph{carried} outwards by the front.

\subsection*{How long it takes}
With $\lambda=1.5$ and $f=1$ the front moves at $c=1.5$ units per second. The
farthest vertex from the new slit, local $(-4,-1.25)$, is the corner $(4,4)$,
at distance $\sqrt{8^2+5.25^2}\approx 9.6$. The front gets there after
$6.4\,$s, and the band behind it after another $\Delta/c=1\,$s. That is the
seven-second camera zoom that follows the opening of the slit.

% =====================================================================
\section{Two clocks: the envelope starts, the phase does not}
% =====================================================================

The crucial design decision is \emph{which} clock each part reads.
\begin{align*}
  \text{phase:}\quad    & \omega t \;\to\; 2\pi f\,(t+t_\text{shift})
                          &&\text{common to all sources}\\
  \text{envelope:}\quad & c\,(t-t_j) &&\text{per source, from its own impact}
\end{align*}
The phase of the late source is \emph{not} restarted at $t_j$. Physically the
second slit is lit by the same plane wave as the first, whose crests reach the
whole wall at once, so the moment the new slit is open it oscillates in step
with the old one. Had its phase been measured from $t_j$ instead, the two
slits would be out of phase by $\omega t_j$, and the central maximum would
move off the symmetry axis.

This is the difference from the other transient the modifier offers,
\node{front=True}, which is a stone dropped into a pond: there the clock is
$t-t_\text{impact}$ for the phase as well, so the wave leaves its source at zero
elongation. The two options are mutually exclusive, and the constructor
raises if both are given.

\node{TimeShift} only enters the phase. It is what lets this scene continue
the wave of \node{towards\_slit} without a jump (see
\node{\_slit\_time\_shift}); the impacts and the fronts stay on the scene time
of the current scene.

% =====================================================================
\section{The node setup}
% =====================================================================

The tree is three frames. Only the first two know about the late source; the
geometry frame is the same as for any other wave.

\subsection{Control frame}

\begin{center}
\begin{tabular}{@{}lll@{}}
\toprule
node (label) & type & role\\
\midrule
\node{Clock}      & Scene Time $\to$ Seconds & $t$, no keyframes needed\\
\node{TimeShift}  & Value & added to $t$ in the phase only\\
\node{Wavelength}, \node{Frequency} & Value & $\lambda$, $f$; also give $c=\lambda f$\\
\node{Amplitude}  & Value & $A$\\
\node{Source0}, \node{Source1} & Vector & $\mathbf c_0$, $\mathbf c_1$ (animatable)\\
\node{FrontWidth} & Value & $\Delta$, built only if some source has an impact\\
\node{Impact1}    & Value & $t_1$; one \node{Impact<j>} per late source\\
\node{Intensity}, \node{IntensityGain} & Value & only with \node{intensity=True}\\
\bottomrule
\end{tabular}
\end{center}

\node{Impact1} is an ordinary value node, so a scene can still move the moment
of opening after the tree has been built:
\verb|ibpy.get_geometry_node_from_modifier(wave2, "Impact1")|.

\subsection{Wave frame}

All of the arithmetic is one \node{make\_function} group named
\node{Elongation}. Its auxiliary expressions, in reverse Polish notation and in
the order in which they are emitted, are for the two sources of the scene
($j=1$ is the late one):

\begin{center}
\small
\begin{tabular}{@{}lll@{}}
\toprule
name & RPN & meaning\\
\midrule
\node{k}      & \verb|tau,wavelength,/|                       & $2\pi/\lambda$\\
\node{wt}     & \verb|time,time_shift,+,frequency,*,tau,*|    & common phase $\omega(t+t_\text{shift})$\\
\node{amp}    & \verb|amplitude,pi,*,wavelength,sqrt,/|       & $A'$\\
\node{r0}, \node{r1} & \verb|pos,c1,sub,length|               & $r_j$\\
\node{x0}, \node{x1} & \verb|r1,0.2,max,k,*|                  & $k\max(r_j,r_\text{src})$\\
\node{reach1} & \verb|wavelength,frequency,*,time,impact1,-,*| & $c\,(t-t_1)$\\
\node{gate1}  & \verb|reach1,r1,-,edge,/,0,max,1,min|         & $\hat x$\\
\node{env1}   & \verb|gate1,gate1,*,3,2,gate1,*,-,*|          & $S(\hat x)$\\
\node{cj$_j$}, \node{sj$_j$} & \verb|amp,x1,j0,*,env1,*| (and \verb|y0|) & gated quadratures\\
\node{w$_j$}  & \verb|cj1,wt,cos,*,sj1,wt,sin,*,+|            & $w_j$ of eq.~\eqref{eq:steady}, gated\\
\node{u}      & \verb|w0,w1,+|                                & elongation\\
\node{C}, \node{S} & \verb|cj0,cj1,+|, \verb|sj0,sj1,+|       & for the intensity\\
\bottomrule
\end{tabular}
\end{center}

The three entries \node{reach1}, \node{gate1}, \node{env1} are produced by
\node{wave\_front\_gate} in \node{geometry\_nodes/nodes.py}, called as
\begin{verbatim}
wave_front_gate("r1", "1", "time,impact1,-")
\end{verbatim}
--- the third argument is exactly where the two clocks part: the gate gets
\verb|time,impact1,-|, while \node{wt} keeps \verb|time,time_shift,+|.
Source~0 has no impact, so it gets no gate and its \node{cj0}, \node{sj0}
carry no \verb|env0| factor.

Two details of the order matter:
\begin{itemize}[leftmargin=*]
  \item The gate entries are inserted into \node{aux\_functions}
        \emph{before} the \node{w$_j$} that multiplies by them;
        \node{make\_function} emits aux expressions in insertion order, and a
        name defined later does not exist yet.
  \item With \node{intensity=True} the gate multiplies the two quadratures
        $C_j=A'J_0\,\mathrm{env}_j$ and $S_j=A'Y_0\,\mathrm{env}_j$ separately,
        not just their combination. So the time-averaged intensity
        $C^2+S^2$, which the \node{Intensity} dial fades the surface to, is
        gated in the same way and grows from the new slit just like the wave.
\end{itemize}

The group has two outputs: \node{elongation} $=u\,(1-I)$ and \node{value}
$=u(1-I)+g\,I\,(C^2+S^2)$, with $I$ the \node{Intensity} dial and $g$ the
\node{IntensityGain}. At $I=0$, the whole of this scene until its last
seconds, both are $u$.

\subsection{Geometry frame}

\[
  \text{Grid}\to\text{store \node{UVMap}}\to\text{store \node{result}}
  \to\text{store \node{resultHeight}}\to\text{store \node{amplitude}}
  \to\text{Set Position}\to\text{Shade Smooth}\to\text{Set Material}
\]
The elongation is stored as an attribute while the grid is still flat --- a
field stored after the lift would measure $\sqrt{r^2+u^2}$ instead of $r$ ---
and \node{Set Position} reads it back through a \node{Named Attribute}.

The colour needs no separate gating. The \node{"function"} material
(\node{function\_texture} in \node{appearance/textures.py}) paints each point
by its \node{result}/\node{amplitude} attribute, i.e.\ by the number the tree
computed, gate included. This is why the scene uses it: the alternative
\node{"interference"} material recomputes the field in the shader and knows
about \node{front}, but not about \node{source\_impacts}, so it would paint the
finished pattern on a surface that has not yet heard of the second slit.

% =====================================================================
\section{A look, not a solution}
% =====================================================================

Equation~\eqref{eq:env} is not what the wave equation does. The true field of
a source switched on at $t_1$ is the retarded convolution with the 2D Green's
function $\propto \Theta(c\tau-r)/\sqrt{c^2\tau^2-r^2}$, which has a sharp
front at $r=c(t-t_1)$ but, unlike in three dimensions, a wake behind it: it
approaches the steady state only slowly, with a tail decaying like a power of
$t$.

What the envelope gets right are the two things an audience actually reads:
nothing changes at distance $r$ from the new slit before $t_1+r/c$, and behind
the front the pattern is the stationary one. A slower, exact relaxation would
spend most of the shot on a transient nobody is meant to look at.

% =====================================================================
\section{Checking it without Blender}
% =====================================================================

Every field built into the tree also exists in numpy.
\node{WaveVisualizationModifier.elongation\_numpy(points, seconds)} evaluates
eq.~\eqref{eq:steady} with \node{scipy.special.j0}/\node{y0} and applies the
same gate per source in \node{source\_impacts}, so the build-up can be
plotted at, say, $t=9,\,11,\,13,\,15\,$s and compared with frames of the
render. (It ignores \node{Intensity}; at $I=0$ that is the whole of what the
tree computes.)

\end{document}
